\noindent\textbf{Scope and editorial date.}
This chapter covers all 28 requested catalogue records and all 94 supplied
source pages, with status assessed on 10 November 2026. It indexes 305
mathematical occurrences, not 305 independent discoveries. The packet is
not a perfectly delimited academic year. Its record envelopes overlap;
record 422 is a presence/heading rather than a second Kirzhnits lecture;
the omitted catalogue record 437 survives on pages 753--754 and is treated
as an explicitly attributed precursor to 438. Records 445, 451 and 457
contain 42 occurrences from 1978--79. Their requested identifiers
are retained, with the actual notebook dates explained. The date field
preserves the supplied catalogue date; it must be read with the provenance
notice, especially for Manin's 21 November 1977 and Fuchs's 3 April 1978
and 5 March 1979 lectures. P000705 retains an unresolved April/September
date conflict. One truncated attachment, P000722, was recovered
from MCCME and checked against its intact neighbors and the uploaded prefix.

\begin{center}
\begin{tabular}{@{}lrrrrrrr@{}}
\toprule
Class & A & B & C & D & E & F & Total\\
\midrule
Occurrences & 55 & 11 & 4 & 0 & 210 & 25 & 305\\
\bottomrule
\end{tabular}
\end{center}
Here A means near-contemporary publication by the speaker or Gelfand;
B, later development by others; C, a precisely delimited partial result;
D, an undeveloped identifiable problem; E, prior established material;
and F, insufficiently identifiable content. These are historical fates,
not correctness ratings. A published thread can contain a false displayed
formula, and an F entry is not evidence of an orphan research idea.

\paragraph{The main lines.}
The strongest connected sequences concern local algebra models for
Harish--Chandra modules, self-dual gauge fields and their finite-dimensional
data, and spectral methods for nonlinear evolution equations. The local
representation-theory announcements match
\citet[Theorems~1--2]{bernsteingelfandgelfand1978local1} and
\citet[Theorem~1]{bernsteingelfandgelfand1978local2}; their small-domain
and regular-central-character qualifications are essential. Manin's talk
precedes the receipt of \citet{drinfeldmanin1978selfdual} by six days.
The instanton construction does not classify every finite-action
Yang--Mills critical point. The current comparison includes
\citet[Theorems~1.1--1.2]{dayapremawaldron2026parabolic}, with their stated
energy and moduli-space scopes kept separate.

Krichever's rational KP theorem, received in September 1977 and published
in 1978, is a close match to the November announcement
\citep[Theorem~3]{krichever1978rational}. Its pole dynamics require a
consistent choice of time and Hamiltonian normalization. Dickey's Poisson
bracket, Faddeev's classical and quantum scattering discussion, and the
Calogero--Degasperis lectures connect infinite-dimensional equations with
finite spectral or particle data. BGG's exterior-algebra account of bundles,
Aomoto's hypergeometric and Ising lectures, Buchstaber's substitution groups,
and Fuchs's vector-field cohomology supply additional substantial lines.
Margulis, Kirzhnits and the counterfeit-coin discussion contain much
established material whose hypotheses still need careful reconstruction.

\paragraph{Questions and development.}
The 13 problem entries include questions answered before the talk,
near-contemporary solutions and surviving restricted questions; they are
not a count of currently open problems. Nonarithmetic lattices in complex
hyperbolic dimension at least four remain open in the located literature
\citep[Introduction]{falbelpasquinelli2023hybrid}; an October 2026 preprint
still reports examples only in dimensions two and three
\citep[Introduction]{yu2026mirror}. The broad Yang--Mills
critical-point classification and the exact scope of soliton-resolution
statements must not be replaced by narrower solved assertions. Buchstaber's
integer-substitution question omits coefficients and cochain topology.
\citet[Theorems~1.2 and~1.5]{babenkobogatyi2008substitutions} resolve its
real finite-jet version and compute a closed-commutator quotient. For an
explicit integral finite-jet interpretation, this chapter proves degree
one is \(\mathbb Z^2\), all degree-two torsion is \((\mathbb Z/2)^2\),
and all products of degree-one classes vanish. These are elementary
consequences and explicit calculations, with no priority claim and no
assertion about the full discrete cohomology ring.

\paragraph{Continuity and critical findings.}
The supplied earlier index connects the spectral-curve material to the
1975 rank-one correspondence and particle systems, gauge theory to the
1975 connection lectures, and the cohomology material to earlier
formal-vector-field discussions. Higher-rank commuting operators and
Hamiltonian-vector-field cohomology are related but different problems;
their old problem IDs are not silently reused. The earlier index is
incomplete as a year-by-year inventory, so a missing match establishes
only that no exact recurrence was identifiable there.

There are 111 critical entries, separating transcription,
translation, ambiguity, provenance and faithful mathematical error.
Among the independently checked cases, thirteen coins is correct for
identity-only detection; the KP coefficients change with convention;
the Calogero equilibrium linearization needs an identity matrix; and two
weight-25 spectral labels coincide, so gap partitions cannot count
distinct Laplace eigenvalues. Exact computations and their outputs are
provided. No D item is inferred from an unsuccessful search. The detailed
report explains the search limits, verification levels and requirements
for presenting these distinctions on the companion website.
\nocite{mccme2026gelfand}
